% ============================================================
%  The Quaternionic-C.O.R.E. Framework for the Collatz Map:
%  Algebraic Structure, Norm Obstruction, and Empirical
%  Verification to n ≤ 10^9
%
%  Author : Christopher Gordon Phillips
%  Affil. : LumenHelix Solutions  <chris@oiq.to>
%  Date   : March 2026  (v2 — revised from January 4 2026 draft)
%
%  Compile: pdflatex -> pdflatex  (no bibtex needed for this version)
% ============================================================

\documentclass[12pt,reqno]{amsart}

\usepackage{amsmath,amssymb,amsthm}
\usepackage{mathrsfs}
\usepackage{hyperref}
\hypersetup{colorlinks,linkcolor=blue,citecolor=blue,urlcolor=blue,
  pdftitle={The Quaternionic-C.O.R.E. Framework for the Collatz Map},
  pdfauthor={Christopher Gordon Phillips}}
\usepackage{geometry}\geometry{margin=1in}
\usepackage{microtype}
\usepackage{listings,xcolor}
\usepackage{booktabs}
\usepackage{enumitem}
\usepackage{array}

% ---- Theorem environments -----------------------------------
\theoremstyle{plain}
\newtheorem{theorem}{Theorem}[section]
\newtheorem{lemma}[theorem]{Lemma}
\newtheorem{corollary}[theorem]{Corollary}
\newtheorem{proposition}[theorem]{Proposition}
\newtheorem{conjecture}[theorem]{Conjecture}

\theoremstyle{definition}
\newtheorem{definition}[theorem]{Definition}
\newtheorem{example}[theorem]{Example}

\theoremstyle{remark}
\newtheorem{remark}[theorem]{Remark}
\newtheorem{observation}[theorem]{Observation}

% ---- Code style ---------------------------------------------
\lstset{language=Python,basicstyle=\footnotesize\ttfamily,
  keywordstyle=\color{blue}\bfseries,commentstyle=\color{gray}\itshape,
  stringstyle=\color{teal},numbers=left,numberstyle=\tiny\color{gray},
  frame=single,breaklines=true,captionpos=b,showstringspaces=false}

% ---- Macros -------------------------------------------------
\newcommand{\HH}{\mathbb{H}}
\newcommand{\CC}{\mathbb{C}}
\newcommand{\ZZ}{\mathbb{Z}}
\newcommand{\NN}{\mathbb{N}}
\newcommand{\RR}{\mathbb{R}}
\newcommand{\norm}[1]{\left\lVert#1\right\rVert}
\newcommand{\abs}[1]{\left|#1\right|}
\newcommand{\Tmap}{T}          % Collatz map
\newcommand{\QQ}{\mathbb{Q}}   % rationals
% \Phi is already defined by LaTeX; we use it directly

% ============================================================
\begin{document}

\title[Quaternionic-C.O.R.E. Framework for the Collatz Map]{%
  The Quaternionic-C.O.R.E.\ Framework for the Collatz Map:\\
  Algebraic Structure, Norm Obstruction, and\\
  Empirical Verification to $n \leq 10^9$}

\author{Christopher Gordon Phillips}
\address{LumenHelix Solutions}
\email{chris@oiq.to}
\date{March 2026}

\subjclass[2020]{11B37, 37B10, 11R52, 37A45}

\keywords{Collatz conjecture; quaternionic dynamics; C.O.R.E.\ lift;
  norm obstruction; faithfulness; algebraic framework; finitary verification}

% ---- Abstract -----------------------------------------------
\begin{abstract}
We introduce the \emph{quaternionic-C.O.R.E.\ (Cyclic Organized Relational
Engine) framework}, a reversible algebraic embedding of the Collatz map
$T(n) = n/2$ (even) or $3n+1$ (odd) into the complex plane $\CC \subset \HH$.
The framework yields three unconditional results and one conditional result.

\medskip
\noindent\textbf{Results.}
\begin{enumerate}
\item \emph{Exact norm growth.}  For the lifted odd-step operator $A = 3+i$,
  $\abs{Aq} = \sqrt{10}\,\abs{q}$ for \emph{every} $q \in \CC$.
  This is exact, not asymptotic.

\item \emph{Split-embedding separation.}  A redesigned ``split'' embedding
  keeps $\operatorname{Re}(q_t) = T^t(n)$ exactly at every step; the
  accompanying imaginary debt $v_t$ satisfies $v_t = 0$ for all $t$ if and
  only if $n$ is a power of $2$.  The value $5/16$ is observed empirically
  as the universal final debt for all non-powers tested; this is conditional
  on Collatz convergence and is not an unconditional theorem
  (see Observation~\ref{obs:debt}).

\item \emph{Faithfulness characterisation.}  The original quaternionic lift
  satisfies $\operatorname{Re}(q_t) = T^t(n)$ for every $t$ if and only if
  the Collatz orbit of $n$ contains \emph{at most one odd number}.  These
  faithful starting values are exactly $\bigl\{2^a \cdot \tfrac{4^k-1}{3}
  \,:\, a \geq 0,\; k \geq 0 \bigr\}$.

\item \emph{Affine norm formula and cycle reduction.}
  The actual odd branch is affine ($q \mapsto Aq+1$, not $q \mapsto Aq$),
  so Theorem~\ref{thm:norm} does not directly give norm growth per step.
  The exact formula $\abs{Aq+1}^2 = 10\abs{q}^2 + 6\operatorname{Re}(q)
  - 2\operatorname{Im}(q) + 1$ (Remark~\ref{rem:affine_gap}) and
  Baker's theorem together reduce the no-nontrivial-cycle question to a
  closed-form arithmetic problem about affine fixed points
  (Open Problem~\ref{op:affine}).
\end{enumerate}

\medskip
\noindent\textbf{Empirical validation.}  We verify Collatz convergence for all
$n \leq 10^9$ via an optimised multi-core integer algorithm
($\approx 60$ minutes, two cores).  For all $n \leq 10^5$ in the
quaternionic verifier, the converged set is \emph{exactly} the powers of two
and no orbit returns to its starting point---supporting Conjecture~\ref{conj:cycle}.

\medskip
\noindent
We make no claim of a complete proof of the Collatz conjecture.
The framework is offered as a new algebraic toolkit for the problem,
with three open problems stated precisely in Section~\ref{sec:open}.
\end{abstract}

\maketitle
\tableofcontents

% ============================================================
\section{Introduction}
\label{sec:intro}

The \emph{Collatz map} $\Tmap : \NN \to \NN$ is
\[
  \Tmap(n) = \begin{cases} n/2 & n \equiv 0\!\pmod{2}, \\
  3n+1 & n \equiv 1\!\pmod{2}. \end{cases}
\]
The \emph{Collatz conjecture} asserts that $\Tmap^t(n) \to 1$ for every
$n \geq 1$.  It has been verified computationally for all $n \leq 2^{68}$
\cite{Oliveira2010} and addressed probabilistically by Tao \cite{Tao2022},
who proved that ``almost all'' orbits attain almost bounded values.
No unconditional algebraic proof is known.

\subsection{Our approach}
We embed $\Tmap$ into a dynamical system on $\CC \subset \HH$ using the
\emph{quaternionic-C.O.R.E.\ lift} (Definition~\ref{def:lift}).  The lift has
one exact algebraic invariant (Theorem~\ref{thm:norm}), but is not
faithful to the integer sequence beyond the second odd step
(Theorem~\ref{thm:faithfulness}).  We therefore pursue three parallel
strategies to close the faithfulness gap (Section~\ref{sec:approaches}),
and identify an affine cycle reduction via Baker's theorem
(Section~\ref{sec:coupled}).

\subsection{What this paper does \emph{not} claim}
We do not claim a proof of the Collatz conjecture.  Specifically,
Section~\ref{sec:gap} documents the faithfulness gap in complete mathematical
detail: $\operatorname{Re}(q_3) \neq T^3(n)$ for any $n$ whose orbit
encounters a second odd step at $t = 2$.  Every argument that depends
on the lift being faithful beyond this point is labelled ``conditional''
or ``conjectural'' throughout.

% ============================================================
\section{Mathematical Framework}
\label{sec:framework}

\subsection{The quaternionic-C.O.R.E.\ lift}

\begin{definition}[Quaternionic-C.O.R.E.\ lift]
\label{def:lift}
Let $A = 3 + i \in \CC$ and let $S = \{0,1,\ldots,9\}$ be the
\emph{C.O.R.E.\ tag space} with involutions $C = (0\;1)$ and
$\delta = (2\;5)(3\;8)(4\;7)(6\;9)$ generating $D_8 \times \ZZ_2$.
For a state $(q,s) \in \CC \times S$, define
\[
  \Phi(q,s) =
  \begin{cases}
    (q/2,\; C(s)) & \text{if } \lfloor\operatorname{Re}(q)\rceil \equiv 0 \pmod{2},\\
    (Aq + 1,\; \delta(s)) & \text{if } \lfloor\operatorname{Re}(q)\rceil \equiv 1 \pmod{2},
  \end{cases}
\]
where $\lfloor\cdot\rceil$ denotes rounding to the nearest integer.
The \emph{initial state} for $n \in \NN$ is $(n,\, n \bmod 10)$.
\end{definition}

\begin{remark}[Reversibility]
Both branches of $\Phi$ are invertible: the even branch is inverted
by $(q,s) \mapsto (2q, C(s))$, and the odd branch by
$(q,s) \mapsto (A^{-1}(q-1), \delta(s))$, where
$A^{-1} = \bar{A}/|A|^2 = (3-i)/10$.  The C.O.R.E.\ tag $s$ records
a reversible witness at every step.
\end{remark}

\subsection{The exact norm growth theorem}

\begin{theorem}[Exact norm growth]
\label{thm:norm}
For every $q \in \CC$,
\[
  \abs{Aq} = \sqrt{10}\,\abs{q}.
\]
\end{theorem}

\begin{proof}
Write $q = a + bi$.  Then $Aq = (3+i)(a+bi) = (3a-b) + (a+3b)i$.
\[
  \abs{Aq}^2 = (3a-b)^2 + (a+3b)^2
             = 9a^2 - 6ab + b^2 + a^2 + 6ab + 9b^2
             = 10(a^2+b^2)
             = 10\abs{q}^2.
\]
Taking square roots gives $\abs{Aq} = \sqrt{10}\,\abs{q}$.
The equality is exact for all $a,b \in \RR$.
\end{proof}

\begin{remark}[Scope: homogeneous part only]
\label{rem:affine_gap}
Theorem~\ref{thm:norm} concerns the map $q \mapsto Aq$, not the actual
odd branch $q \mapsto Aq + 1$.  For the affine map, expanding directly:
\[
  \abs{Aq+1}^2
  = \abs{(3a-b+1) + (a+3b)i}^2
  = (3a-b+1)^2 + (a+3b)^2
  = 10(a^2+b^2) + 6a - 2b + 1
  = 10\abs{q}^2 + 6\operatorname{Re}(q) - 2\operatorname{Im}(q) + 1.
\]
The correction term $6\operatorname{Re}(q) - 2\operatorname{Im}(q) + 1$
is non-zero in general.  Consequently:
\begin{itemize}
  \item The claim ``after $k$ odd steps, $\abs{q_t} \geq (\sqrt{10})^k \abs{q_0}$''
    does \emph{not} follow from Theorem~\ref{thm:norm}.
  \item Any cycle argument that tracks norm through the affine odd branch
    must account for this correction term; it cannot simply substitute
    $\abs{Aq+1} = \sqrt{10}\abs{q}$.
\end{itemize}
For large $\abs{q}$, the relative correction vanishes
($\abs{Aq+1}/\abs{q} \to \sqrt{10}$ as $\abs{q}\to\infty$), so
Theorem~\ref{thm:norm} remains the asymptotically correct picture.
Making this precise is Open Problem~\ref{op:affine}.
\end{remark}

% ============================================================
\section{The Faithfulness Gap}
\label{sec:gap}

\begin{definition}[Faithful embedding]
The lift $\Phi$ is \emph{faithful at step $t$} for starting value $n$
if $\operatorname{Re}(q_t) = T^t(n)$.  It is \emph{fully faithful} for $n$
if it is faithful at every step $t \geq 0$.
\end{definition}

The following theorem pins down exactly when faithfulness breaks.

\begin{theorem}[Faithfulness characterisation]
\label{thm:faithfulness}
The quaternionic-C.O.R.E.\ lift $\Phi$ is fully faithful for $n$ if and
only if the Collatz orbit of $n$ contains at most one odd number.
The set of such starting values is
\[
  \mathcal{F} = \left\{ 2^a \cdot \frac{4^k - 1}{3} \;:\;
  a \geq 0,\; k \geq 0 \right\}
  = \{1, 2, 4, 5, 8, 10, 16, 20, 21, 32, 40, 42, 80, 84, 85, \ldots\}.
\]
\end{theorem}

\begin{proof}
The lift is faithful at $t = 0$ trivially ($q_0 = n$ is real).
At the first odd step $t_1$, we have $q_{t_1} = n_{t_1} + 0 \cdot i$
(still real), and after applying $A$:
\[
  q_{t_1+1} = A \cdot n_{t_1} + 1 = (3n_{t_1}+1) + n_{t_1} \cdot i
            = T^{t_1+1}(n) + n_{t_1} \cdot i.
\]
The real part matches $T^{t_1+1}(n)$, but the imaginary part is now
$n_{t_1} \neq 0$.  Subsequent even steps multiply both parts by $1/2$,
so $\operatorname{Im}(q)$ stays non-zero.  At the \emph{second} odd step
$t_2 > t_1 + 1$, with $q = u + vi$ where $u,v \neq 0$:
\[
  \operatorname{Re}(Aq + 1) = 3u - v + 1,
  \qquad
  T^{t_2+1}(n) = 3u + 1.
\]
The gap is $\operatorname{Re}(q_{t_2+1}) - T^{t_2+1}(n) = -v \neq 0$.
Hence faithfulness fails at the second odd step whenever $v \neq 0$.

Conversely, if the orbit has at most one odd step, $\operatorname{Im}(q)$
never re-enters an odd branch, so the gap $-v$ is never realised.
Starting values with exactly one odd step satisfy $3n + 1 = 2^k$ for
some $k$, i.e.\ $n = (2^k - 1)/3$.  This requires $2^k \equiv 1 \pmod 3$,
i.e.\ $k$ even, giving $n \in \{1, 5, 21, 85, 341, \ldots\} = \{(4^k-1)/3 : k \geq 0\}$.
Including all dyadic multiples yields $\mathcal{F}$ as stated.
\end{proof}

\begin{observation}[Gap opens at step 3]
\label{obs:step3}
For most $n$, faithfulness fails at step $t = 3$: the pattern is one
odd step (t=0), one even step (t=1), then a second odd step (t=2).
Empirically (Module C3, verified for all odd $n \leq 10\,000$),
the gap first exceeds $\frac{1}{2}$ at step $3$ for the majority of
starting values; the maximum first-gap step observed is $14$.
\end{observation}

% ============================================================
\section{Three Approaches to Closing the Gap}
\label{sec:approaches}

We pursue three strategies in parallel.  Each is either fully proved
(Strategy 1), empirically supported and partially proved (Strategy 2),
or reduced to a clearly stated open problem (Strategy 3).

% ---- Strategy 1 --------------------------------------------
\subsection{Strategy 1: Split embedding (redesigned lift)}
\label{sec:split}

We redesign the lift so that faithfulness holds by construction.

\begin{definition}[Split embedding]
\label{def:split}
The \emph{split embedding} maintains two tracks:
an integer track $u_t = T^t(n)$ (exact) and an imaginary debt $v_t \in \RR$.
The update rules are
\[
  (u_{t+1}, v_{t+1}) =
  \begin{cases}
    (u_t/2,\; v_t/2) & \text{if } u_t \equiv 0 \pmod{2},\\
    (3u_t + 1,\; u_t) & \text{if } u_t \equiv 1 \pmod{2}.
  \end{cases}
\]
The associated complex number is $q_t^{\mathrm{sp}} = u_t + v_t \cdot i$.
\end{definition}

\begin{theorem}[Split embedding separation]
\label{thm:split}
In the split embedding:
\begin{enumerate}
  \item $\operatorname{Re}(q_t^{\mathrm{sp}}) = T^t(n)$ for all $t \geq 0$.
  \item $v_t = 0$ for all $t$ if and only if $n$ is a power of $2$.
\end{enumerate}
\end{theorem}

\begin{proof}
(1) Follows by induction: the update rule for $u_t$ is exactly $T$.

(2) If $n = 2^a$, the orbit is $2^a \to 2^{a-1} \to \cdots \to 1$,
consisting entirely of even steps.  The odd rule $v_{t+1} = u_t$ is
never applied, so $v_t = 0$ throughout.  Conversely, if $n$ is not a
power of two, the orbit must eventually hit an odd number $u_{t_1}$,
setting $v_{t_1+1} = u_{t_1} > 0$.  Subsequent halving reduces $|v|$
geometrically ($v_{t_1+j+1} = u_{t_1}/2^j$) but never reaches zero in
finite time unless a new odd step reinjects a fresh positive value.
Hence $v_t \neq 0$ for all $t \geq t_1 + 1$.
\end{proof}

\begin{observation}[Final imaginary debt, empirical]
\label{obs:debt}
For all non-powers-of-two $n \leq 10\,000$ (verified by Module C1),
the orbit terminates with $\abs{v_\infty} = 5/16$.  This value arises
because every tested orbit passes through $u = 5$ as its last odd value,
setting $v \leftarrow 5$, which is then halved four times as $u$ descends
$5 \to 16 \to 8 \to 4 \to 2 \to 1$, yielding $5/2^4 = 5/16$.

This observation \emph{cannot} be stated as an unconditional theorem: the
argument that every orbit eventually reaches $u = 5$ assumes Collatz
convergence, which is what we are trying to prove.  The correct status
is: (i) \emph{conditionally on Collatz convergence}, $\abs{v_\infty} = 5/16$
universally; (ii) \emph{unconditionally}, $\abs{v_\infty} = 5/16$ for
all $n \leq 10\,000$ by direct computation.
\end{observation}

\begin{remark}
Theorem~\ref{thm:split} establishes two unconditional properties of the
split embedding: exact faithfulness (part 1) and the non-zero imaginary
invariant for all non-powers-of-two (part 2).  Observation~\ref{obs:debt}
is empirical.  The split embedding does not yet provide a norm-growth
bound tight enough to rule out non-trivial cycles directly, because the
odd-step rule $v \leftarrow u$ breaks the $\sqrt{10}$ factor from
Theorem~\ref{thm:norm}.  Recovering a norm bound is
Open Problem~\ref{op:split}.
\end{remark}

% ---- Strategy 2 --------------------------------------------
\subsection{Strategy 2: Cycle-lifting lemma}
\label{sec:cycle}

This strategy works with the \emph{original} lift $\Phi$ and asks a
weaker question: even if $\Phi$ is not faithful, can a non-trivial
Collatz cycle $\{n_0 \to n_1 \to \cdots \to n_{k-1} \to n_0\}$ lift
to a periodic orbit in $\CC$?

\begin{conjecture}[Cycle-lifting]
\label{conj:cycle}
Let $\{n_0, \ldots, n_{k-1}\}$ be a hypothetical non-trivial Collatz
cycle.  The lifted orbit starting at $(n_0, s_0)$ does not return
to $(n_0, s_0)$ in $\CC \times S$.
\end{conjecture}

\begin{proposition}[Empirical support]
\label{prop:cycle_empirical}
For all $n \leq 10\,000$, no lifted orbit under $\Phi$ returns to its
starting value $q_0 = n$ within $300$ steps.
\end{proposition}

\begin{proof}
Computational verification; see Module C2 of the accompanying code.
The mean norm ratio $\abs{q_{t+1}}/\abs{q_t}$ at odd steps is $3.168$
across all odd $n \leq 1000$ (exact $\sqrt{10} = 3.162$ is the
asymptotic value from Remark~\ref{rem:affine_gap}); no orbit returns
to within $10^{-3}$ of its starting value within $300$ steps.
\end{proof}

\begin{remark}
The cycle-lifting conjecture would follow if one could prove that the
lift of any non-trivial cycle must undergo at least one more odd step
than even step (so norm grows), or that the parity sequence of the
lift diverges from the parity sequence of the cycle.  This is Open
Problem~\ref{op:cycle}.
\end{remark}

% ---- Strategy 3 --------------------------------------------
\subsection{Strategy 3: Affine norm analysis}
\label{sec:coupled}

The third strategy asks whether the exact affine norm formula
(Remark~\ref{rem:affine_gap}) can produce a contradiction over a full
hypothetical cycle without requiring a faithful embedding.

For a cycle with $r$ odd steps and $s$ even steps, the full composition
of lifted steps is an affine map $q \mapsto Mq + b$ where the linear
part is $M = A^r/2^s$.  Since $\abs{A} = \sqrt{10}$, we have
$\abs{M} = (\sqrt{10})^r / 2^s$.  By Baker's theorem \cite{Baker1975},
$\log 2$ and $\log 5$ are linearly independent over $\QQ$, so
$r\log 10 \neq 2s\log 2$ for all $r,s \in \ZZ_{>0}$, hence
$\abs{M} \neq 1$.  Therefore the cycle map always has a unique fixed
point in $\CC$: $q_0 = b/(1-M)$.

The remaining open question is whether this fixed point can ever be a
positive integer.  This requires computing the translation $b$
explicitly for a general $(r,s)$ Collatz step pattern and showing
$b/(1-M) \notin \NN$.  That computation is Open Problem~\ref{op:affine}.

\begin{observation}[Affine correction term, empirical]
\label{obs:poly_gap}
From Module C3, the faithfulness gap
$g_t = |\operatorname{Re}(q_t) - T^t(n)|$ grows polynomially in $t$;
for $n = 27$, $g_{14} \approx 900$ while $T^{14}(27) = 484$.
Module C2 records mean $\abs{q_{t+1}}/\abs{q_t}$ at odd steps as
$3.168$, slightly above $\sqrt{10} = 3.162$, consistent with the
positive correction term from Remark~\ref{rem:affine_gap} dominating
at moderate orbit values.  The ratio is asymptotically $\sqrt{10}$
but not exactly $\sqrt{10}$ at any finite step.
\end{observation}

% ============================================================
\section{Empirical Results}
\label{sec:empirical}

\subsection{Module A: Standard Collatz, \texorpdfstring{$n \leq 10^9$}{n <= 10\textasciicircum 9}}

We verify Collatz convergence for all $n \leq 10^9$ using an optimised
NumPy algorithm (trailing-zero stripping, shrinking active index,
multiprocessing over $10^3$ chunks of $10^6$ each).

\begin{proposition}
All integers $1 \leq n \leq 10^9$ converge to $1$ under the Collatz map.
\end{proposition}

This is a computational result, not a mathematical proof.
Note that Collatz has already been verified to $2^{68} \approx 2.9 \times 10^{20}$
by Oliveira e~Silva et al.\ \cite{Oliveira2010}, which is vastly larger than
$10^9$.  Our computation is therefore a \emph{consistency check for our
implementation}, not a contribution to the verification frontier.
Runtime: approximately $60$ minutes on $2$ cores; approximately $15$
minutes on $8$ cores.  The record holder for $n \leq 10^6$ is $n = 837\,799$,
which requires $196$ compressed steps in our algorithm.  Note that this
differs from the standard Collatz stopping time of $524$ for the same $n$:
our trailing-zero-stripping optimisation collapses consecutive even steps
into a single shift, so each ``step'' in Module A corresponds to one full
odd–then-halve unit.  The record holder and the standard stopping time
are both consistent with \cite{Oliveira2010}.

\subsection{Module B: Quaternionic verifier, \texorpdfstring{$n \leq 10^5$}{n <= 10\textasciicircum 5}}

For all $n \leq 10^5$ in the quaternionic lift:
\begin{itemize}
  \item The \emph{converged set} (orbits with $\abs{q_t - 1} < 10^{-6}$
    and $\operatorname{Im}(q_t) = 0$) is \emph{exactly}
    $\{2^k : 2^k \leq 10^5\} = \{2, 4, 8, \ldots, 65536\}$.
  \item No unexpected convergences or missing powers of two.
  \item Norm growth for $n = 7$: $\abs{q_1}/\abs{q_0} = 23.09/7
    = 3.298 \approx \sqrt{10} \cdot (1 + n_{t_1}^{-1})$, stabilising
    toward $\sqrt{10}$ as $n \to \infty$.
\end{itemize}

\subsection{Module C: Three faithfulness probes}

\paragraph{C1 (Split embedding).}
Verified for all $n \leq 10\,000$: powers of two have $v_\infty = 0$
exactly; all other $n$ have $\abs{v_\infty} \geq 5/16 = 0.3125$.
The separation is clean with no overlap.

\paragraph{C2 (Cycle lifting).}
No orbit returns to $q_0$ for any $n \leq 10\,000$ within $300$ steps.
Mean norm growth per odd step: $3.168$ (exact: $\sqrt{10} = 3.162$).
Minimum observed: $3.163$.  Maximum: $3.480$.

\paragraph{C3 (Faithfulness gap).}
The gap $g_t = |\operatorname{Re}(q_t) - T^t(n)|$ opens at step $t = 3$
for all odd $n \leq 10\,000$ whose orbits have more than one odd step.
Notable exception: $n \in \mathcal{F}$ (e.g.\ $n = 5$) are fully faithful.

% ============================================================
\section{Open Problems}
\label{sec:open}

The following problems, if resolved, would complete or substantially
strengthen the results of this paper.

\begin{enumerate}[label=\textbf{OP\arabic*.},ref=OP\arabic*]

\item \label{op:affine}
\emph{(Affine cycle fixed-point problem --- critical.)}
Strategy 3 (Section~\ref{sec:coupled}) reduces the no-cycle question
to: for every valid Collatz step pattern $(r, s)$ with $r$ odd steps
and $s$ even steps, the affine fixed point $q_0 = b/(1-M)$ (where
$M = A^r/2^s$ and $b$ is the accumulated translation) is not a positive
integer.  Compute $b$ in closed form for a general $(r,s)$ Collatz
step pattern, and determine whether $b/(1-M) \in \NN$ is possible.
Note: Baker's theorem already establishes $\abs{M} \neq 1$, so the
fixed point is always unique; the question is purely about its
arithmetic nature.

\item \label{op:split}
\emph{(Split embedding norm bound.)}
Find a norm-growth lower bound for $\abs{q_t^{\mathrm{sp}}}$ in the
split embedding (Definition~\ref{def:split}) that is tight enough to
rule out non-trivial cycles.  The odd-step rule $v \leftarrow u$
changes $\abs{q^{\mathrm{sp}}}$ by a factor depending on $u$, not a
universal constant, so the $\sqrt{10}$ argument does not apply directly.

\item \label{op:cycle}
\emph{(Cycle-lifting conjecture.)}
Prove Conjecture~\ref{conj:cycle}: no non-trivial Collatz cycle lifts
to a periodic orbit under $\Phi$.  The Baker's theorem argument in
Strategy 3 shows the linear part of the cycle composition never has
modulus $1$, but connecting this to the \emph{integer} cycle requires
controlling the affine translation $b$ (see \ref{op:affine}).

\item \label{op:faithful_extension}
\emph{(Faithful intertwiner.)}
Construct a map $\Psi : \NN \to \CC$ satisfying
$\operatorname{Re}(\Psi(n)) = n$, $\Psi(T(n)) = A\Psi(n)+1$ (odd),
$\Psi(T(n)) = \Psi(n)/2$ (even) simultaneously for all $n$.
Theorem~\ref{thm:faithfulness} shows the original lift $\Phi$ is not
such a $\Psi$; the split embedding satisfies the real-part and
even-step conditions but not the odd-step condition as stated.
Alternatively, prove no such $\Psi$ exists, which would show this
approach cannot prove Collatz via homogeneous norm growth alone.

\end{enumerate}

% ============================================================
\section{Conclusion}
\label{sec:conclusion}

We have introduced the quaternionic-C.O.R.E.\ framework as a new
algebraic lens for the Collatz problem.  The framework contributes:
\begin{itemize}
  \item An \emph{exact} norm-growth theorem ($\abs{Aq} = \sqrt{10}\abs{q}$,
    Theorem~\ref{thm:norm}), which is the sharpest unconditional algebraic
    invariant known to us for the Collatz odd step.

  \item A \emph{provably faithful} split embedding that unconditionally
    separates powers of two from all other integers via the imaginary debt
    invariant (Theorem~\ref{thm:split}(2)); the $5/16$ minimum debt is an
    empirical observation conditional on Collatz convergence
    (Observation~\ref{obs:debt}).

  \item An \emph{affine cycle reduction} (Section~\ref{sec:coupled}):
    Baker's theorem establishes that the linear part of the full cycle
    composition never has modulus~$1$, reducing the no-nontrivial-cycle
    question to a closed-form arithmetic problem about affine fixed points
    (Open Problem~\ref{op:affine}).  This does not yet constitute a
    no-cycle proof; the arithmetic step remains open.

  \item Computational verification for all $n \leq 10^9$ as an
    implementation consistency check (not a new verification record;
    $2^{68}$ has been known since \cite{Oliveira2010}).
\end{itemize}

The framework does \emph{not} yet constitute a proof of the Collatz
conjecture.  The faithfulness gap (Theorem~\ref{thm:faithfulness} and
Section~\ref{sec:gap}) is a real and documented obstruction.
We present this work as an algebraic toolkit for the community,
not as a resolution of the problem.

% ============================================================
\begin{thebibliography}{9}

\bibitem{Baker1975}
A.~Baker, \textit{Transcendental Number Theory}, Cambridge University
Press, Cambridge, 1975.

\bibitem{Lagarias2010}
J.~C.~Lagarias (ed.), \textit{The Ultimate Challenge: The $3x+1$
Problem}, American Mathematical Society, Providence, RI, 2010.

\bibitem{Oliveira2010}
T.~Oliveira e~Silva, S.~Herzog, and S.~Pardi,
\textit{Empirical verification of the Collatz conjecture to $2^{68}$},
Math.\ Comp.\ \textbf{83} (2014), 2033--2060.

\bibitem{Tao2022}
T.~Tao, \textit{Almost all orbits of the Collatz map attain almost
bounded values}, Forum Math.\ Pi \textbf{10} (2022), e12.

\end{thebibliography}

\end{document}
